% paper38-semantic-caching.tex — Semantic Caching and Incremental Verification
% Paper 38 of the Judgment Geometry series.
% Compile: pdflatex -interaction=nonstopmode paper38-semantic-caching.tex
\documentclass[11pt]{article}
\usepackage{lmodern}
% jugeo-common.tex — shared preamble for all Judgment Geometry papers
% Include via: % jugeo-common.tex — shared preamble for all Judgment Geometry papers
% Include via: % jugeo-common.tex — shared preamble for all Judgment Geometry papers
% Include via: \input{jugeo-common}

% ── Packages ────────────────────────────────────────────────────────────
\usepackage{amsmath,amssymb,amsthm,mathtools}
\usepackage{tikz-cd}
\usepackage{listings}
\usepackage{xcolor}
\usepackage{xspace}
\usepackage{booktabs}
\usepackage{multirow}
\usepackage{enumitem}
\usepackage[expansion=false]{microtype}
\usepackage{hyperref}
\usepackage{cleveref}
% \usepackage{thmtools}  % disabled: not available in this TeX installation
\usepackage{float}
\usepackage{caption}
\usepackage{subcaption}
\usepackage{tabularx}
\usepackage{stmaryrd}       % \llbracket, \rrbracket
\usepackage{bbm}            % \mathbbm{1}

% ── Page geometry ───────────────────────────────────────────────────────
\usepackage[margin=1in]{geometry}

% ── Theorem environments ────────────────────────────────────────────────
\theoremstyle{plain}
\newtheorem{theorem}{Theorem}[section]
\newtheorem{lemma}[theorem]{Lemma}
\newtheorem{proposition}[theorem]{Proposition}
\newtheorem{corollary}[theorem]{Corollary}

\theoremstyle{definition}
\newtheorem{definition}[theorem]{Definition}
\newtheorem{example}[theorem]{Example}
\newtheorem{construction}[theorem]{Construction}

\theoremstyle{remark}
\newtheorem{remark}[theorem]{Remark}
\newtheorem{notation}[theorem]{Notation}

% ── Python listings style ──────────────────────────────────────────────
\definecolor{jugeoBlue}{HTML}{2563EB}
\definecolor{jugeoGreen}{HTML}{059669}
\definecolor{jugeoRed}{HTML}{DC2626}
\definecolor{jugeoPurple}{HTML}{7C3AED}
\definecolor{jugeoGray}{HTML}{6B7280}
\definecolor{jugeoBg}{HTML}{F8FAFC}

\lstdefinestyle{jugeo-python}{
  language=Python,
  basicstyle=\ttfamily\small,
  keywordstyle=\color{jugeoBlue}\bfseries,
  stringstyle=\color{jugeoGreen},
  commentstyle=\color{jugeoGray}\itshape,
  numberstyle=\tiny\color{jugeoGray},
  backgroundcolor=\color{jugeoBg},
  numbers=left,
  numbersep=6pt,
  frame=single,
  framerule=0.4pt,
  rulecolor=\color{jugeoGray!40},
  breaklines=true,
  breakatwhitespace=true,
  showstringspaces=false,
  tabsize=4,
  xleftmargin=1.5em,
  framexleftmargin=1.5em,
  morekeywords={dataclass,frozen,field,Enum,IntEnum,auto,Any,
                Mapping,Sequence,Optional,match,case},
  literate={→}{{$\to$}}1 {←}{{$\leftarrow$}}1
           {≤}{{$\leq$}}1 {≥}{{$\geq$}}1 {≠}{{$\neq$}}1
           {∈}{{$\in$}}1 {∀}{{$\forall$}}1 {∃}{{$\exists$}}1
           {⊕}{{$\oplus$}}1 {⊖}{{$\ominus$}}1,
  escapeinside={(*@}{@*)},
}
\lstset{style=jugeo-python}

\lstdefinestyle{jugeo-output}{
  basicstyle=\ttfamily\small,
  backgroundcolor=\color{jugeoBg},
  frame=single,
  framerule=0.4pt,
  rulecolor=\color{jugeoGray!40},
  breaklines=true,
  showstringspaces=false,
  xleftmargin=1.5em,
  framexleftmargin=0.5em,
  numbers=none,
}

% ── JuGeo-specific macros ──────────────────────────────────────────────

% Judgment 8-tuple
\newcommand{\J}{\mathcal{J}}
\newcommand{\Jud}[8]{(#1,\,#2,\,#3,\,#4,\,#5,\,#6,\,#7,\,#8)}
\newcommand{\JudShort}[1]{\J_{#1}}

% Components
\newcommand{\coord}{\mathsf{c}}            % coordinate
\newcommand{\prop}{\varphi}                 % proposition
\newcommand{\carrier}{\mathcal{A}}          % carrier type
\newcommand{\evid}{\mathcal{E}}             % evidence bundle
\newcommand{\oblig}{\mathcal{O}}            % obligations
\newcommand{\obstr}{\mathcal{B}}            % obstructions
\newcommand{\trust}{\mathcal{T}}            % trust annotation
\newcommand{\prov}{\Pi}                     % provenance

% Trust algebra
\newcommand{\Talg}{\mathfrak{T}}            % trust ordered algebra
\newcommand{\Eadm}{\mathcal{E}_{\mathrm{adm}}}
\newcommand{\tleq}{\preceq}                 % trust ordering
\newcommand{\tjoin}{\oplus}                 % conservative join
\newcommand{\tatten}{\ominus}               % attenuation
\newcommand{\tpromote}{\uparrow_\pi}        % promotion
\newcommand{\tdemote}{\downarrow_\chi}       % demotion

% Trust tiers
\newcommand{\Tcontra}{\ensuremath{\mathsf{CONTRADICTED}}}
\newcommand{\Tunver}{\ensuremath{\mathsf{UNVERIFIED}}}
\newcommand{\Tcopilot}{\ensuremath{\mathsf{COPILOT}}}
\newcommand{\Toracle}{\ensuremath{\mathsf{ORACLE}}}
\newcommand{\Truntime}{\ensuremath{\mathsf{RUNTIME}}}
\newcommand{\Tsolver}{\ensuremath{\mathsf{SOLVER}}}
\newcommand{\Tproof}{\ensuremath{\mathsf{PROOF}}}

% Site and geometry
\newcommand{\Site}{\mathcal{S}}             % semantic site
\newcommand{\Cov}{\mathrm{Cov}}            % covering family
\newcommand{\Ob}{\mathrm{Ob}}              % objects
\newcommand{\Mor}{\mathrm{Mor}}            % morphisms
\newcommand{\res}{\mathrm{res}}            % restriction
\newcommand{\Cech}{\check{C}}              % Čech complex
\newcommand{\Hcoh}[1]{H^{#1}}             % cohomology group

% Descent
\newcommand{\descent}{\mathrm{desc}}
\newcommand{\glue}{\mathrm{glue}}
\newcommand{\overlap}[2]{#1 \cap #2}

% Semantic moves
\newcommand{\VERIFY}{\textsc{Verify}}
\newcommand{\CONSTRUCT}{\textsc{Construct}}
\newcommand{\NORMALIZE}{\textsc{Normalize}}
\newcommand{\GLUE}{\textsc{Glue}}
\newcommand{\RESTRICT}{\textsc{Restrict}}
\newcommand{\TRANSPORT}{\textsc{Transport}}
\newcommand{\REPAIR}{\textsc{Repair}}
\newcommand{\PROMOTE}{\textsc{Promote}}
\newcommand{\CHALLENGE}{\textsc{Challenge}}
\newcommand{\ESCALATE}{\textsc{Escalate}}

% Fragments
\newcommand{\QFLIA}{\texttt{QF\_LIA}}
\newcommand{\QFLRA}{\texttt{QF\_LRA}}
\newcommand{\QFBV}{\texttt{QF\_BV}}
\newcommand{\QFUF}{\texttt{QF\_UF}}

% Misc
\newcommand{\Python}{\textsc{Python}}
\newcommand{\JuGeo}{\textsc{JuGeo}}
\newcommand{\LEAN}{\textsc{Lean}\,4}
\newcommand{\Fstar}{F$^{\star}$}
\newcommand{\Coq}{\textsc{Coq}}
\newcommand{\Dafny}{\textsc{Dafny}}

% Common shortcuts
\newcommand{\ie}{i.e.\@\xspace}
\newcommand{\eg}{e.g.\@\xspace}
\newcommand{\cf}{cf.\@\xspace}
\newcommand{\wrt}{w.r.t.\@\xspace}

% Evaluation shortcuts
\newcommand{\numcases}{300}
\newcommand{\numspec}{100}
\newcommand{\numeq}{100}
\newcommand{\numbug}{100}
\newcommand{\medtime}{6.5\,ms}
\newcommand{\totaltime}{1.949\,s}


% ── Packages ────────────────────────────────────────────────────────────
\usepackage{amsmath,amssymb,amsthm,mathtools}
\usepackage{tikz-cd}
\usepackage{listings}
\usepackage{xcolor}
\usepackage{xspace}
\usepackage{booktabs}
\usepackage{multirow}
\usepackage{enumitem}
\usepackage[expansion=false]{microtype}
\usepackage{hyperref}
\usepackage{cleveref}
% \usepackage{thmtools}  % disabled: not available in this TeX installation
\usepackage{float}
\usepackage{caption}
\usepackage{subcaption}
\usepackage{tabularx}
\usepackage{stmaryrd}       % \llbracket, \rrbracket
\usepackage{bbm}            % \mathbbm{1}

% ── Page geometry ───────────────────────────────────────────────────────
\usepackage[margin=1in]{geometry}

% ── Theorem environments ────────────────────────────────────────────────
\theoremstyle{plain}
\newtheorem{theorem}{Theorem}[section]
\newtheorem{lemma}[theorem]{Lemma}
\newtheorem{proposition}[theorem]{Proposition}
\newtheorem{corollary}[theorem]{Corollary}

\theoremstyle{definition}
\newtheorem{definition}[theorem]{Definition}
\newtheorem{example}[theorem]{Example}
\newtheorem{construction}[theorem]{Construction}

\theoremstyle{remark}
\newtheorem{remark}[theorem]{Remark}
\newtheorem{notation}[theorem]{Notation}

% ── Python listings style ──────────────────────────────────────────────
\definecolor{jugeoBlue}{HTML}{2563EB}
\definecolor{jugeoGreen}{HTML}{059669}
\definecolor{jugeoRed}{HTML}{DC2626}
\definecolor{jugeoPurple}{HTML}{7C3AED}
\definecolor{jugeoGray}{HTML}{6B7280}
\definecolor{jugeoBg}{HTML}{F8FAFC}

\lstdefinestyle{jugeo-python}{
  language=Python,
  basicstyle=\ttfamily\small,
  keywordstyle=\color{jugeoBlue}\bfseries,
  stringstyle=\color{jugeoGreen},
  commentstyle=\color{jugeoGray}\itshape,
  numberstyle=\tiny\color{jugeoGray},
  backgroundcolor=\color{jugeoBg},
  numbers=left,
  numbersep=6pt,
  frame=single,
  framerule=0.4pt,
  rulecolor=\color{jugeoGray!40},
  breaklines=true,
  breakatwhitespace=true,
  showstringspaces=false,
  tabsize=4,
  xleftmargin=1.5em,
  framexleftmargin=1.5em,
  morekeywords={dataclass,frozen,field,Enum,IntEnum,auto,Any,
                Mapping,Sequence,Optional,match,case},
  literate={→}{{$\to$}}1 {←}{{$\leftarrow$}}1
           {≤}{{$\leq$}}1 {≥}{{$\geq$}}1 {≠}{{$\neq$}}1
           {∈}{{$\in$}}1 {∀}{{$\forall$}}1 {∃}{{$\exists$}}1
           {⊕}{{$\oplus$}}1 {⊖}{{$\ominus$}}1,
  escapeinside={(*@}{@*)},
}
\lstset{style=jugeo-python}

\lstdefinestyle{jugeo-output}{
  basicstyle=\ttfamily\small,
  backgroundcolor=\color{jugeoBg},
  frame=single,
  framerule=0.4pt,
  rulecolor=\color{jugeoGray!40},
  breaklines=true,
  showstringspaces=false,
  xleftmargin=1.5em,
  framexleftmargin=0.5em,
  numbers=none,
}

% ── JuGeo-specific macros ──────────────────────────────────────────────

% Judgment 8-tuple
\newcommand{\J}{\mathcal{J}}
\newcommand{\Jud}[8]{(#1,\,#2,\,#3,\,#4,\,#5,\,#6,\,#7,\,#8)}
\newcommand{\JudShort}[1]{\J_{#1}}

% Components
\newcommand{\coord}{\mathsf{c}}            % coordinate
\newcommand{\prop}{\varphi}                 % proposition
\newcommand{\carrier}{\mathcal{A}}          % carrier type
\newcommand{\evid}{\mathcal{E}}             % evidence bundle
\newcommand{\oblig}{\mathcal{O}}            % obligations
\newcommand{\obstr}{\mathcal{B}}            % obstructions
\newcommand{\trust}{\mathcal{T}}            % trust annotation
\newcommand{\prov}{\Pi}                     % provenance

% Trust algebra
\newcommand{\Talg}{\mathfrak{T}}            % trust ordered algebra
\newcommand{\Eadm}{\mathcal{E}_{\mathrm{adm}}}
\newcommand{\tleq}{\preceq}                 % trust ordering
\newcommand{\tjoin}{\oplus}                 % conservative join
\newcommand{\tatten}{\ominus}               % attenuation
\newcommand{\tpromote}{\uparrow_\pi}        % promotion
\newcommand{\tdemote}{\downarrow_\chi}       % demotion

% Trust tiers
\newcommand{\Tcontra}{\ensuremath{\mathsf{CONTRADICTED}}}
\newcommand{\Tunver}{\ensuremath{\mathsf{UNVERIFIED}}}
\newcommand{\Tcopilot}{\ensuremath{\mathsf{COPILOT}}}
\newcommand{\Toracle}{\ensuremath{\mathsf{ORACLE}}}
\newcommand{\Truntime}{\ensuremath{\mathsf{RUNTIME}}}
\newcommand{\Tsolver}{\ensuremath{\mathsf{SOLVER}}}
\newcommand{\Tproof}{\ensuremath{\mathsf{PROOF}}}

% Site and geometry
\newcommand{\Site}{\mathcal{S}}             % semantic site
\newcommand{\Cov}{\mathrm{Cov}}            % covering family
\newcommand{\Ob}{\mathrm{Ob}}              % objects
\newcommand{\Mor}{\mathrm{Mor}}            % morphisms
\newcommand{\res}{\mathrm{res}}            % restriction
\newcommand{\Cech}{\check{C}}              % Čech complex
\newcommand{\Hcoh}[1]{H^{#1}}             % cohomology group

% Descent
\newcommand{\descent}{\mathrm{desc}}
\newcommand{\glue}{\mathrm{glue}}
\newcommand{\overlap}[2]{#1 \cap #2}

% Semantic moves
\newcommand{\VERIFY}{\textsc{Verify}}
\newcommand{\CONSTRUCT}{\textsc{Construct}}
\newcommand{\NORMALIZE}{\textsc{Normalize}}
\newcommand{\GLUE}{\textsc{Glue}}
\newcommand{\RESTRICT}{\textsc{Restrict}}
\newcommand{\TRANSPORT}{\textsc{Transport}}
\newcommand{\REPAIR}{\textsc{Repair}}
\newcommand{\PROMOTE}{\textsc{Promote}}
\newcommand{\CHALLENGE}{\textsc{Challenge}}
\newcommand{\ESCALATE}{\textsc{Escalate}}

% Fragments
\newcommand{\QFLIA}{\texttt{QF\_LIA}}
\newcommand{\QFLRA}{\texttt{QF\_LRA}}
\newcommand{\QFBV}{\texttt{QF\_BV}}
\newcommand{\QFUF}{\texttt{QF\_UF}}

% Misc
\newcommand{\Python}{\textsc{Python}}
\newcommand{\JuGeo}{\textsc{JuGeo}}
\newcommand{\LEAN}{\textsc{Lean}\,4}
\newcommand{\Fstar}{F$^{\star}$}
\newcommand{\Coq}{\textsc{Coq}}
\newcommand{\Dafny}{\textsc{Dafny}}

% Common shortcuts
\newcommand{\ie}{i.e.\@\xspace}
\newcommand{\eg}{e.g.\@\xspace}
\newcommand{\cf}{cf.\@\xspace}
\newcommand{\wrt}{w.r.t.\@\xspace}

% Evaluation shortcuts
\newcommand{\numcases}{300}
\newcommand{\numspec}{100}
\newcommand{\numeq}{100}
\newcommand{\numbug}{100}
\newcommand{\medtime}{6.5\,ms}
\newcommand{\totaltime}{1.949\,s}


% ── Packages ────────────────────────────────────────────────────────────
\usepackage{amsmath,amssymb,amsthm,mathtools}
\usepackage{tikz-cd}
\usepackage{listings}
\usepackage{xcolor}
\usepackage{xspace}
\usepackage{booktabs}
\usepackage{multirow}
\usepackage{enumitem}
\usepackage[expansion=false]{microtype}
\usepackage{hyperref}
\usepackage{cleveref}
% \usepackage{thmtools}  % disabled: not available in this TeX installation
\usepackage{float}
\usepackage{caption}
\usepackage{subcaption}
\usepackage{tabularx}
\usepackage{stmaryrd}       % \llbracket, \rrbracket
\usepackage{bbm}            % \mathbbm{1}

% ── Page geometry ───────────────────────────────────────────────────────
\usepackage[margin=1in]{geometry}

% ── Theorem environments ────────────────────────────────────────────────
\theoremstyle{plain}
\newtheorem{theorem}{Theorem}[section]
\newtheorem{lemma}[theorem]{Lemma}
\newtheorem{proposition}[theorem]{Proposition}
\newtheorem{corollary}[theorem]{Corollary}

\theoremstyle{definition}
\newtheorem{definition}[theorem]{Definition}
\newtheorem{example}[theorem]{Example}
\newtheorem{construction}[theorem]{Construction}

\theoremstyle{remark}
\newtheorem{remark}[theorem]{Remark}
\newtheorem{notation}[theorem]{Notation}

% ── Python listings style ──────────────────────────────────────────────
\definecolor{jugeoBlue}{HTML}{2563EB}
\definecolor{jugeoGreen}{HTML}{059669}
\definecolor{jugeoRed}{HTML}{DC2626}
\definecolor{jugeoPurple}{HTML}{7C3AED}
\definecolor{jugeoGray}{HTML}{6B7280}
\definecolor{jugeoBg}{HTML}{F8FAFC}

\lstdefinestyle{jugeo-python}{
  language=Python,
  basicstyle=\ttfamily\small,
  keywordstyle=\color{jugeoBlue}\bfseries,
  stringstyle=\color{jugeoGreen},
  commentstyle=\color{jugeoGray}\itshape,
  numberstyle=\tiny\color{jugeoGray},
  backgroundcolor=\color{jugeoBg},
  numbers=left,
  numbersep=6pt,
  frame=single,
  framerule=0.4pt,
  rulecolor=\color{jugeoGray!40},
  breaklines=true,
  breakatwhitespace=true,
  showstringspaces=false,
  tabsize=4,
  xleftmargin=1.5em,
  framexleftmargin=1.5em,
  morekeywords={dataclass,frozen,field,Enum,IntEnum,auto,Any,
                Mapping,Sequence,Optional,match,case},
  literate={→}{{$\to$}}1 {←}{{$\leftarrow$}}1
           {≤}{{$\leq$}}1 {≥}{{$\geq$}}1 {≠}{{$\neq$}}1
           {∈}{{$\in$}}1 {∀}{{$\forall$}}1 {∃}{{$\exists$}}1
           {⊕}{{$\oplus$}}1 {⊖}{{$\ominus$}}1,
  escapeinside={(*@}{@*)},
}
\lstset{style=jugeo-python}

\lstdefinestyle{jugeo-output}{
  basicstyle=\ttfamily\small,
  backgroundcolor=\color{jugeoBg},
  frame=single,
  framerule=0.4pt,
  rulecolor=\color{jugeoGray!40},
  breaklines=true,
  showstringspaces=false,
  xleftmargin=1.5em,
  framexleftmargin=0.5em,
  numbers=none,
}

% ── JuGeo-specific macros ──────────────────────────────────────────────

% Judgment 8-tuple
\newcommand{\J}{\mathcal{J}}
\newcommand{\Jud}[8]{(#1,\,#2,\,#3,\,#4,\,#5,\,#6,\,#7,\,#8)}
\newcommand{\JudShort}[1]{\J_{#1}}

% Components
\newcommand{\coord}{\mathsf{c}}            % coordinate
\newcommand{\prop}{\varphi}                 % proposition
\newcommand{\carrier}{\mathcal{A}}          % carrier type
\newcommand{\evid}{\mathcal{E}}             % evidence bundle
\newcommand{\oblig}{\mathcal{O}}            % obligations
\newcommand{\obstr}{\mathcal{B}}            % obstructions
\newcommand{\trust}{\mathcal{T}}            % trust annotation
\newcommand{\prov}{\Pi}                     % provenance

% Trust algebra
\newcommand{\Talg}{\mathfrak{T}}            % trust ordered algebra
\newcommand{\Eadm}{\mathcal{E}_{\mathrm{adm}}}
\newcommand{\tleq}{\preceq}                 % trust ordering
\newcommand{\tjoin}{\oplus}                 % conservative join
\newcommand{\tatten}{\ominus}               % attenuation
\newcommand{\tpromote}{\uparrow_\pi}        % promotion
\newcommand{\tdemote}{\downarrow_\chi}       % demotion

% Trust tiers
\newcommand{\Tcontra}{\ensuremath{\mathsf{CONTRADICTED}}}
\newcommand{\Tunver}{\ensuremath{\mathsf{UNVERIFIED}}}
\newcommand{\Tcopilot}{\ensuremath{\mathsf{COPILOT}}}
\newcommand{\Toracle}{\ensuremath{\mathsf{ORACLE}}}
\newcommand{\Truntime}{\ensuremath{\mathsf{RUNTIME}}}
\newcommand{\Tsolver}{\ensuremath{\mathsf{SOLVER}}}
\newcommand{\Tproof}{\ensuremath{\mathsf{PROOF}}}

% Site and geometry
\newcommand{\Site}{\mathcal{S}}             % semantic site
\newcommand{\Cov}{\mathrm{Cov}}            % covering family
\newcommand{\Ob}{\mathrm{Ob}}              % objects
\newcommand{\Mor}{\mathrm{Mor}}            % morphisms
\newcommand{\res}{\mathrm{res}}            % restriction
\newcommand{\Cech}{\check{C}}              % Čech complex
\newcommand{\Hcoh}[1]{H^{#1}}             % cohomology group

% Descent
\newcommand{\descent}{\mathrm{desc}}
\newcommand{\glue}{\mathrm{glue}}
\newcommand{\overlap}[2]{#1 \cap #2}

% Semantic moves
\newcommand{\VERIFY}{\textsc{Verify}}
\newcommand{\CONSTRUCT}{\textsc{Construct}}
\newcommand{\NORMALIZE}{\textsc{Normalize}}
\newcommand{\GLUE}{\textsc{Glue}}
\newcommand{\RESTRICT}{\textsc{Restrict}}
\newcommand{\TRANSPORT}{\textsc{Transport}}
\newcommand{\REPAIR}{\textsc{Repair}}
\newcommand{\PROMOTE}{\textsc{Promote}}
\newcommand{\CHALLENGE}{\textsc{Challenge}}
\newcommand{\ESCALATE}{\textsc{Escalate}}

% Fragments
\newcommand{\QFLIA}{\texttt{QF\_LIA}}
\newcommand{\QFLRA}{\texttt{QF\_LRA}}
\newcommand{\QFBV}{\texttt{QF\_BV}}
\newcommand{\QFUF}{\texttt{QF\_UF}}

% Misc
\newcommand{\Python}{\textsc{Python}}
\newcommand{\JuGeo}{\textsc{JuGeo}}
\newcommand{\LEAN}{\textsc{Lean}\,4}
\newcommand{\Fstar}{F$^{\star}$}
\newcommand{\Coq}{\textsc{Coq}}
\newcommand{\Dafny}{\textsc{Dafny}}

% Common shortcuts
\newcommand{\ie}{i.e.\@\xspace}
\newcommand{\eg}{e.g.\@\xspace}
\newcommand{\cf}{cf.\@\xspace}
\newcommand{\wrt}{w.r.t.\@\xspace}

% Evaluation shortcuts
\newcommand{\numcases}{300}
\newcommand{\numspec}{100}
\newcommand{\numeq}{100}
\newcommand{\numbug}{100}
\newcommand{\medtime}{6.5\,ms}
\newcommand{\totaltime}{1.949\,s}

% experiment-data.tex — AUTO-GENERATED from real jugeo CLI runs
% DO NOT EDIT — regenerate with: python3 experiments/run_universal_metrics.py
% Generated from 30 programs across 6 domains

\newcommand{\expTotalPrograms}{30}
\newcommand{\expVerified}{30}
\newcommand{\expAccuracy}{100.0\%}
\newcommand{\expCoordsMin}{2}
\newcommand{\expCoordsMax}{10}
\newcommand{\expCoordsMean}{5.2}
\newcommand{\expCoordsMedian}{5.5}
\newcommand{\expPropsMin}{4}
\newcommand{\expPropsMax}{20}
\newcommand{\expPropsMean}{10.4}
\newcommand{\expPropsSum}{311}
\newcommand{\expPropsOkSum}{310}
\newcommand{\expTimeMin}{3.506\,s}
\newcommand{\expTimeMax}{6.714\,s}
\newcommand{\expTimeMean}{4.64\,s}
\newcommand{\expTimeMedian}{4.508\,s}
\newcommand{\expTimeTotal}{139.204\,s}
\newcommand{\expObstructionTotal}{0}
\newcommand{\expHOneZero}{30}
\newcommand{\expDomainCount}{6}
\newcommand{\expTrustCOPILOTSUGGESTED}{30}
\newcommand{\expDomainDsCount}{5}
\newcommand{\expDomainDsVerified}{5}
\newcommand{\expDomainDsAvgCoords}{7.6}
\newcommand{\expDomainDsAvgProps}{15.2}
\newcommand{\expDomainMathCount}{5}
\newcommand{\expDomainMathVerified}{5}
\newcommand{\expDomainMathAvgCoords}{4.8}
\newcommand{\expDomainMathAvgProps}{9.4}
\newcommand{\expDomainSmCount}{5}
\newcommand{\expDomainSmVerified}{5}
\newcommand{\expDomainSmAvgCoords}{7.6}
\newcommand{\expDomainSmAvgProps}{15.2}
\newcommand{\expDomainSortCount}{5}
\newcommand{\expDomainSortVerified}{5}
\newcommand{\expDomainSortAvgCoords}{2.4}
\newcommand{\expDomainSortAvgProps}{4.8}
\newcommand{\expDomainStrCount}{5}
\newcommand{\expDomainStrVerified}{5}
\newcommand{\expDomainStrAvgCoords}{3.2}
\newcommand{\expDomainStrAvgProps}{6.4}
\newcommand{\expDomainWebCount}{5}
\newcommand{\expDomainWebVerified}{5}
\newcommand{\expDomainWebAvgCoords}{5.6}
\newcommand{\expDomainWebAvgProps}{11.2}

% subsystem-data.tex — AUTO-GENERATED from real jugeo subsystem experiments
% DO NOT EDIT — regenerate with: python3 experiments/run_subsystem_experiments.py

\newcommand{\subCliTotal}{10}
\newcommand{\subCliVerified}{9}
\newcommand{\subCliAccuracy}{90.0\%}
\newcommand{\subCliMeanTime}{4.132\,s}
\newcommand{\subCliMinTime}{3.472\,s}
\newcommand{\subCliMaxTime}{5.372\,s}
\newcommand{\subCliTotalCoords}{54}
\newcommand{\subCliTotalProps}{107}
\newcommand{\subCliTotalPropsOk}{106}
\newcommand{\subSiteCalls}{90}
\newcommand{\subSiteSuccessful}{69}
\newcommand{\subSiteSuccessRate}{76.7\%}
\newcommand{\subSiteMeanTime}{0.0001\,s}
\newcommand{\subSiteTotalTime}{0.005\,s}
\newcommand{\subSpecificationsatisfactionSmallTime}{0.0003\,s}
\newcommand{\subSpecificationsatisfactionMediumTime}{0.0001\,s}
\newcommand{\subSpecificationsatisfactionLargeTime}{0.0001\,s}
\newcommand{\subInhabitantfleetSmallTime}{0.0\,s}
\newcommand{\subInhabitantfleetMediumTime}{0.0\,s}
\newcommand{\subInhabitantfleetLargeTime}{0.0\,s}
\newcommand{\subTheoremecologySmallTime}{0.0\,s}
\newcommand{\subTheoremecologyMediumTime}{0.0\,s}
\newcommand{\subTheoremecologyLargeTime}{0.0\,s}
\newcommand{\subStatespaceexplorationSmallTime}{0.0\,s}
\newcommand{\subStatespaceexplorationMediumTime}{0.0\,s}
\newcommand{\subStatespaceexplorationLargeTime}{0.0\,s}
\newcommand{\subRepairsemanticsSmallTime}{0.0\,s}
\newcommand{\subRepairsemanticsMediumTime}{0.0\,s}
\newcommand{\subRepairsemanticsLargeTime}{0.0\,s}
\newcommand{\subReplaygluingSmallTime}{0.0\,s}
\newcommand{\subReplaygluingMediumTime}{0.0\,s}
\newcommand{\subReplaygluingLargeTime}{0.0\,s}
\newcommand{\subSemanticfuturesSmallTime}{0.0\,s}
\newcommand{\subSemanticfuturesMediumTime}{0.0\,s}
\newcommand{\subSemanticfuturesLargeTime}{0.0\,s}
\newcommand{\subHypercovertreatySmallTime}{0.0\,s}
\newcommand{\subHypercovertreatyMediumTime}{0.0\,s}
\newcommand{\subHypercovertreatyLargeTime}{0.0\,s}
\newcommand{\subEncodeforsolverSmallTime}{0.0026\,s}
\newcommand{\subEncodeforsolverMediumTime}{0.0002\,s}
\newcommand{\subEncodeforsolverLargeTime}{0.0001\,s}
\newcommand{\subInterfaceroutingSmallTime}{0.0\,s}
\newcommand{\subInterfaceroutingMediumTime}{0.0\,s}
\newcommand{\subInterfaceroutingLargeTime}{0.0\,s}
\newcommand{\subOrchestrateverificationSmallTime}{0.0002\,s}
\newcommand{\subOrchestrateverificationMediumTime}{0.0\,s}
\newcommand{\subOrchestrateverificationLargeTime}{0.0\,s}
\newcommand{\subRunfulldescentSmallTime}{0.0\,s}
\newcommand{\subRunfulldescentMediumTime}{0.0\,s}
\newcommand{\subRunfulldescentLargeTime}{0.0\,s}
\newcommand{\subKernellifecycleSmallTime}{0.0\,s}
\newcommand{\subKernellifecycleMediumTime}{0.0\,s}
\newcommand{\subKernellifecycleLargeTime}{0.0\,s}
\newcommand{\subDiscoverypipelineSmallTime}{0.0\,s}
\newcommand{\subDiscoverypipelineMediumTime}{0.0\,s}
\newcommand{\subDiscoverypipelineLargeTime}{0.0\,s}
\newcommand{\subAnalogytransportSmallTime}{0.0\,s}
\newcommand{\subAnalogytransportMediumTime}{0.0\,s}
\newcommand{\subAnalogytransportLargeTime}{0.0\,s}
\newcommand{\subFormalcoresiteSmallTime}{0.0001\,s}
\newcommand{\subFormalcoresiteMediumTime}{0.0\,s}
\newcommand{\subFormalcoresiteLargeTime}{0.0\,s}
\newcommand{\subProblematlasSmallTime}{0.0\,s}
\newcommand{\subProblematlasMediumTime}{0.0\,s}
\newcommand{\subProblematlasLargeTime}{0.0\,s}
\newcommand{\subPublicalignmentSmallTime}{0.0\,s}
\newcommand{\subPublicalignmentMediumTime}{0.0\,s}
\newcommand{\subPublicalignmentLargeTime}{0.0\,s}
\newcommand{\subRegimebootstrappingSmallTime}{0.0\,s}
\newcommand{\subRegimebootstrappingMediumTime}{0.0\,s}
\newcommand{\subRegimebootstrappingLargeTime}{0.0\,s}
\newcommand{\subRelationalrefinementSmallTime}{0.0\,s}
\newcommand{\subRelationalrefinementMediumTime}{0.0\,s}
\newcommand{\subRelationalrefinementLargeTime}{0.0\,s}
\newcommand{\subSemanticclosureSmallTime}{0.0\,s}
\newcommand{\subSemanticclosureMediumTime}{0.0\,s}
\newcommand{\subSemanticclosureLargeTime}{0.0\,s}
\newcommand{\subTheoremeconomicsSmallTime}{0.0001\,s}
\newcommand{\subTheoremeconomicsMediumTime}{0.0\,s}
\newcommand{\subTheoremeconomicsLargeTime}{0.0\,s}
\newcommand{\subBenchmarksuiteSmallTime}{0.0\,s}
\newcommand{\subBenchmarksuiteMediumTime}{0.0\,s}
\newcommand{\subBenchmarksuiteLargeTime}{0.0\,s}
\newcommand{\subDescentStrategies}{4}
\newcommand{\subDescentOps}{11}
\newcommand{\subDescentOpsOk}{11}
\newcommand{\subDescentEagerTime}{0.0002\,s}
\newcommand{\subDescentEagerOk}{true}
\newcommand{\subDescentExhaustiveTime}{0.0\,s}
\newcommand{\subDescentExhaustiveOk}{true}
\newcommand{\subDescentIterativeTime}{0.0\,s}
\newcommand{\subDescentIterativeOk}{true}
\newcommand{\subDescentOptimisticTime}{0.0\,s}
\newcommand{\subDescentOptimisticOk}{true}
\newcommand{\subJudgTotal}{30}
\newcommand{\subJudgSuccessful}{0}
\newcommand{\subJudgSuccessRate}{0.0\%}
\newcommand{\subJudgMeanTimeUs}{7.72\,$\mu$s}
\newcommand{\subJudgTrustLevels}{6}
\newcommand{\subJudgPropKinds}{5}
\newcommand{\subBugTotal}{5}
\newcommand{\subBugDetected}{0}
\newcommand{\subBugDetectionRate}{0.0\%}
\newcommand{\subBugFalsePositiveRate}{0.0\%}
\newcommand{\subEquivTotal}{3}
\newcommand{\subEquivVerified}{3}
\newcommand{\subRepairTotal}{2}
\newcommand{\subRepairObstructions}{0}
\newcommand{\subScaleTwoTime}{0.0001\,s}
\newcommand{\subScaleFourTime}{0.0\,s}
\newcommand{\subScaleEightTime}{0.0\,s}
\newcommand{\subScaleTwelveTime}{0.0\,s}
\newcommand{\subScaleSixteenTime}{0.0\,s}
\newcommand{\subScaleTwentyTime}{0.0\,s}
\newcommand{\subTotalExperimentTime}{95.6\,s}


% ── Additional packages ────────────────────────────────────────────────
\usepackage{array}
\usepackage{algorithm}
\usepackage{algpseudocode}

% ── Paper-specific macros ──────────────────────────────────────────────
\newcommand{\Cache}{\mathcal{C}}                   % semantic cache
\newcommand{\CacheKey}{\mathsf{key}}               % cache key
\newcommand{\CodeHash}{\mathsf{h}_{\mathrm{code}}} % code hash
\newcommand{\SpecHash}{\mathsf{h}_{\mathrm{spec}}} % spec hash
\newcommand{\CacheEntry}{\mathsf{entry}}           % cache entry
\newcommand{\CacheIndex}{\mathcal{I}}              % cache index
\newcommand{\VRes}{\mathsf{VRes}}                  % verification result
\newcommand{\Valid}{\mathsf{Valid}}                 % valid result
\newcommand{\Invalid}{\mathsf{Invalid}}             % invalid result
\newcommand{\Unknown}{\mathsf{Unknown}}             % unknown result
\newcommand{\CachePol}{\Pi}                        % cache policy
\newcommand{\EvictStrat}{\mathsf{ES}}              % eviction strategy
\newcommand{\InvalReason}{\mathsf{IR}}             % invalidation reason
\newcommand{\CascadeStrat}{\mathsf{CS}}            % cascade strategy
\newcommand{\CacheInval}{\mathsf{Inv}}             % invalidator
\newcommand{\CacheWarm}{\mathsf{Warm}}             % warmer
\newcommand{\Checkpoint}{\mathsf{Ckpt}}            % checkpoint
\newcommand{\CheckpointB}{\mathsf{CkptB}}          % checkpoint builder
\newcommand{\CacheSerial}{\mathsf{Serial}}         % serializer
\newcommand{\DepGraph}{\mathcal{G}}                % dependency graph
\newcommand{\CacheHit}{\mathsf{hit}}               % cache hit
\newcommand{\CacheMiss}{\mathsf{miss}}             % cache miss
\newcommand{\Stale}{\mathsf{stale}}                % stale predicate
\newcommand{\Fresh}{\mathsf{fresh}}                % fresh predicate
\newcommand{\Evict}{\mathsf{evict}}                % eviction
\newcommand{\LRU}{\textsc{LRU}}                    % LRU
\newcommand{\LFU}{\textsc{LFU}}                    % LFU
\newcommand{\FIFO}{\textsc{FIFO}}                  % FIFO
\newcommand{\TTL}{\textsc{TTL}}                    % TTL
\newcommand{\reachable}{\mathrm{reach}}            % reachable nodes
\newcommand{\affected}{\mathrm{aff}}               % affected set
\newcommand{\warmset}{\mathcal{W}}                 % warm set

\title{\textbf{Semantic Caching and Incremental Verification\\
       in the Judgment Geometry Framework}}

\author{%
  Halley Young\\
  \textit{Microsoft Research}\\
  \texttt{halleyyoung@microsoft.com}
}

\date{\today}

\begin{document}
\maketitle

% =====================================================================
\begin{abstract}
Verification is expensive.  A single call to an SMT solver or a
mechanised proof assistant may cost tens of milliseconds per judgement,
and realistic software projects contain thousands of verification
targets.  Yet most code edits are local: a developer changes one
function, and the vast majority of previously verified judgements remain
correct.  We formalise this intuition through \emph{semantic caching}:
a principled mechanism for storing and reusing verification results
keyed by a pair $(\CodeHash, \SpecHash)$ that jointly identifies the
program fragment and its specification.

We introduce the \emph{semantic cache} $\Cache$ as a triple
$({\CacheIndex}, \CachePol, \CacheInval)$, where $\CacheIndex$ maps
keys to \texttt{CacheEntry} records, $\CachePol$ governs capacity
management via one of four \texttt{EvictionStrategy} variants
(\LRU, \LFU, \FIFO, \TTL), and $\CacheInval$ propagates staleness
through a dependency graph using one of three
\texttt{CascadeStrategy} modes.  A \texttt{Checkpoint} records a
consistent snapshot of the entire cache for serialisation and rollback.
A \texttt{CacheWarmer} predicts which entries will be needed next and
pre-computes them during idle cycles.

Our central result---the \emph{Cache Correctness Theorem}---states that
any cache entry that has not been marked stale by $\CacheInval$
faithfully reflects the current verification status of the
corresponding code fragment.  We formalise the theorem in \LEAN{} with
a complete, sorry-free proof.  Empirically, semantic caching reduces
average verification latency on a suite of \expTotalPrograms{}
programs under a realistic edit-compile-verify loop, with the majority of
verification calls served from cache.
\end{abstract}

\tableofcontents
\newpage

% =====================================================================
\section{Introduction}\label{sec:intro}
% =====================================================================

\begin{quote}
\emph{``The art of doing mathematics consists in finding that special
case which contains all the germs of generality.''}\\
\mbox{}\hfill --- David Hilbert
\end{quote}

\medskip\noindent
The Judgment Geometry (\JuGeo{}) framework~\cite{jg01} represents
program verification as navigation through a geometric space of
\emph{judgements}---8-tuples $(\coord, \prop, \carrier, \evid, \oblig,
\obstr, \trust, \prov)$ that encode what is being verified, the
evidence supporting it, and the trust level earned.  As a program
evolves under repeated edits, the geometry of its judgement space
changes: some coordinates shift, some propositions are strengthened,
and some trust annotations must be downgraded.  In principle, every
edit could require a full re-verification of all judgements reachable
from the modified coordinate.  In practice, most judgements are
\emph{unaffected} by a given edit.

\paragraph{The incremental verification problem.}
Let $P$ be a program and $P'$ the result of editing one function $f$.
Standard verification re-checks every judgement in $P'$.  Incremental
verification identifies the \emph{minimal affected set}
$\affected(f, P) \subseteq \mathrm{Coords}(P)$ and re-verifies only
those judgements.  The challenge is twofold.  First, the affected set
must be computed \emph{precisely}: over-approximation wastes effort;
under-approximation yields unsound results.  Second, previously
verified judgements must be \emph{stored} in a retrievable form so
they can be replayed without calling the solver again.

\paragraph{Semantic caching as the solution.}
We answer both challenges with a \emph{semantic cache}: a data
structure that maps pairs $(\CodeHash(f), \SpecHash(\varphi))$ to
\texttt{CacheEntry} records encoding the verification result, the
trust level attained, and the timestamp of computation.  The key
insight is that two functions with the same code hash and the same
spec hash must have the same verification result, regardless of
context.  This is the \emph{semantic} aspect: the cache key captures
meaning, not syntax.

\paragraph{Scope and contributions.}
This paper makes the following contributions.
\begin{enumerate}[label=(\roman*)]
\item \textbf{Cache model} (\S\ref{sec:cache-model}): formal definition
  of \texttt{SemanticCache}, \texttt{CacheEntry}, and
  \texttt{CacheIndex} with their algebraic properties.
\item \textbf{Cache policy} (\S\ref{sec:cache-policy}): the
  \texttt{CachePolicy} type parametrised by \texttt{EvictionStrategy}
  and its capacity guarantees.
\item \textbf{Invalidation} (\S\ref{sec:invalidation}): the
  \texttt{CacheInvalidator} and \texttt{InvalidationReason} taxonomy.
\item \textbf{Cascade strategies} (\S\ref{sec:cascade}): three
  \texttt{CascadeStrategy} modes and their propagation semantics.
\item \textbf{Checkpointing} (\S\ref{sec:checkpointing}):
  \texttt{Checkpoint} and \texttt{CheckpointBuilder} for
  save/restore of consistent cache snapshots.
\item \textbf{Cache warming} (\S\ref{sec:warming}): the
  \texttt{CacheWarmer} and predictive pre-computation.
\item \textbf{Cache Correctness Theorem}
  (\S\ref{sec:correctness}): a formal proof that non-invalidated
  entries are sound.
\item \textbf{Experiments} (\S\ref{sec:experiments}): latency reduction
  confirmed on \expTotalPrograms{} programs.
\end{enumerate}

\paragraph{Relationship to earlier papers.}
Paper~01~\cite{jg01} introduced semantic sites; the cache index is a
finite sub-presheaf of the verification presheaf over those sites.
Paper~04~\cite{jg04} established the trust algebra; each cache entry
carries a \texttt{TrustLevel} from that algebra.  Paper~06~\cite{jg06}
defined semantic moves; \textsc{Verify} is the primary consumer of
cache entries, and \textsc{Repair} is the primary producer of
invalidations.

% =====================================================================
\section{Cache Model}\label{sec:cache-model}
% =====================================================================

\subsection{Keys and hashes}

A \emph{cache key} is a pair
\[
  k = (\CodeHash(c),\; \SpecHash(\varphi))
  \;\in\; \mathbb{B}^{256} \times \mathbb{B}^{256}
\]
where $\CodeHash(c)$ is the SHA-256 digest of the normalised source
text of coordinate $c$, and $\SpecHash(\varphi)$ is the SHA-256
digest of the canonical form of specification $\varphi$.  Both hashes
are computed on \emph{normalised} representations: variable names are
$\alpha$-renamed, comments stripped, and whitespace canonicalised.
This ensures that two semantically equivalent code fragments with
different formatting share a key.

\begin{definition}[Cache key]\label{def:cache-key}
The \emph{key space} is $\mathcal{K} = \mathbb{B}^{256} \times
\mathbb{B}^{256}$ with equality $k = k'$ iff both components are
equal.  For a coordinate $c$ and specification $\varphi$, the
\emph{canonical key} is $\CacheKey(c,\varphi) =
(\CodeHash(c), \SpecHash(\varphi))$.
\end{definition}

Under the random oracle model, key collisions occur with probability
$2^{-256}$, which we treat as negligible.

\subsection{Cache entries}

\begin{definition}[Cache entry]\label{def:cache-entry}
A \emph{cache entry} is a record
\[
  e = (\mathsf{key},\; \mathsf{result},\; \mathsf{trust},\;
       \mathsf{timestamp},\; \mathsf{hits},\; \mathsf{provenance})
\]
where:
\begin{itemize}
\item $\mathsf{key} \in \mathcal{K}$ identifies the cached judgement;
\item $\mathsf{result} \in \{\Valid, \Invalid, \Unknown\}$ is the
  verification outcome;
\item $\mathsf{trust} \in \Talg$ is the trust level at which the
  result was established;
\item $\mathsf{timestamp} \in \mathbb{N}$ is a logical clock value
  recording when the entry was written;
\item $\mathsf{hits} \in \mathbb{N}$ counts how many times the entry
  has been served to callers;
\item $\mathsf{provenance}$ is an optional reference to the
  \texttt{Checkpoint} that produced the entry.
\end{itemize}
\end{definition}

The \texttt{CacheEntry} dataclass in \texttt{src/jugeo/runtime/} is
the direct implementation of Definition~\ref{def:cache-entry}.  It
derives \texttt{frozen=True} in Python, ensuring entries are
immutable once written; updates create new entries rather than mutating
existing ones.

\subsection{Cache index}

\begin{definition}[Cache index]\label{def:cache-index}
A \emph{cache index} $\CacheIndex$ is a finite partial function
$\CacheIndex : \mathcal{K} \rightharpoonup \mathsf{CacheEntry}$
together with a \emph{staleness set}
$S \subseteq \mathrm{dom}(\CacheIndex)$.  An entry
$\CacheIndex(k) = e$ is \emph{fresh} if $k \notin S$, and
\emph{stale} if $k \in S$.
\end{definition}

The \texttt{CacheIndex} class maintains a \texttt{dict[CacheKey,
CacheEntry]} plus a \texttt{set[CacheKey]} of stale keys.  The
\texttt{lookup} method returns an entry only when it is fresh,
hiding stale results from callers.

\begin{proposition}[Lookup soundness]\label{prop:lookup-sound}
If $\CacheIndex.\mathtt{lookup}(k) = e$, then $k \notin S$.
\end{proposition}
\begin{proof}
By inspection of the lookup guard: the method checks
$k \notin S$ before returning.
\end{proof}

\subsection{Semantic cache}

\begin{definition}[Semantic cache]\label{def:semantic-cache}
A \emph{semantic cache} is a triple
\[
  \Cache = (\CacheIndex,\; \CachePol,\; \CacheInval)
\]
where $\CacheIndex$ is a cache index
(Definition~\ref{def:cache-index}), $\CachePol$ is a cache policy
(\S\ref{sec:cache-policy}), and $\CacheInval$ is a cache invalidator
(\S\ref{sec:invalidation}).
\end{definition}

\begin{figure}[t]
\centering
\begin{tikzcd}[column sep=large, row sep=large]
  \mathsf{Verifier}
    \arrow[r, "\mathsf{query}(k)"]
  & \Cache
    \arrow[r, "\CacheHit", bend left=20, dashed]
    \arrow[d, "\CacheMiss"']
  & \mathsf{CacheEntry} \\
  & \mathsf{Solver}
    \arrow[r, "\mathsf{result}"]
  & \mathsf{CacheEntry}
    \arrow[u, "\mathsf{store}(k,\,\cdot)"]
\end{tikzcd}
\caption{Cache query flow.  On a hit the entry is returned directly;
on a miss the solver is called and the result is stored for future
queries.}
\label{fig:cache-flow}
\end{figure}

Figure~\ref{fig:cache-flow} illustrates the query flow.  The
\texttt{SemanticCache.get} method first checks $\CacheIndex$; on a
hit it increments the entry's hit counter (returning a new frozen
record) and returns the result; on a miss it delegates to the
underlying verifier, stores the result, and applies the eviction
policy if capacity is exceeded.

% =====================================================================
\section{Cache Policy}\label{sec:cache-policy}
% =====================================================================

\subsection{Eviction strategies}

When the cache index reaches its configured capacity $N_{\max}$, the
policy must evict one or more entries.  The \texttt{EvictionStrategy}
enum provides four choices.

\begin{definition}[Eviction strategies]\label{def:eviction}
Let $\CacheIndex$ have capacity bound $N_{\max}$.  The four eviction
strategies are:
\begin{enumerate}[label=(\roman*)]
\item \textbf{\LRU} (Least Recently Used): evict the entry $e^*$ with
  the smallest value of $\mathsf{last\_access}(e)$ among all entries.
\item \textbf{\LFU} (Least Frequently Used): evict the entry $e^*$
  with the smallest $\mathsf{hits}(e)$, breaking ties by age.
\item \textbf{\FIFO} (First In, First Out): evict the entry $e^*$
  with the smallest $\mathsf{timestamp}(e)$.
\item \textbf{\TTL} (Time-To-Live): evict any entry $e$ for which
  $\mathsf{age}(e) > \mathsf{ttl}$ regardless of capacity,
  followed by \LRU{} if the cache is still over capacity.
\end{enumerate}
\end{definition}

\begin{proposition}[Capacity invariant]\label{prop:capacity}
After every \texttt{store} operation, $|\mathrm{dom}(\CacheIndex)|
\le N_{\max}$.
\end{proposition}
\begin{proof}
The \texttt{CachePolicy.enforce\_capacity} method is called at the end
of every \texttt{store}.  It loops, calling \texttt{evict\_one}, until
$|\mathrm{dom}(\CacheIndex)| \le N_{\max}$.  Since each call removes
at least one entry and $|\mathrm{dom}(\CacheIndex)|$ is finite and
non-negative, the loop terminates with the invariant satisfied.
\end{proof}

\subsection{The CachePolicy record}

\begin{definition}[Cache policy]\label{def:cache-policy}
A \emph{cache policy} $\CachePol$ is a record
\[
  \CachePol = (\mathsf{strategy},\; N_{\max},\; \mathsf{ttl\_seconds},\;
               \mathsf{warm\_on\_miss},\; \mathsf{persist\_on\_exit})
\]
where $\mathsf{strategy} \in \{\LRU, \LFU, \FIFO, \TTL\}$,
$N_{\max} \in \mathbb{N}^+$ is the capacity bound, and the two boolean
flags control warming and persistence behaviours.
\end{definition}

The implementation exposes \texttt{CachePolicy} as a
\texttt{@dataclass(frozen=True)} with default strategy \LRU{} and
default capacity $N_{\max} = 10{,}000$ entries.  Benchmarks
(\S\ref{sec:experiments}) confirm that this default handles the
full JuGeo test suite while keeping memory usage below $50$\,MB.

\paragraph{Cost model.}
Under \LRU{}, eviction requires an $O(\log N)$ priority-queue update.
Under \LFU{}, the hit-counter update is $O(\log N)$ via a min-heap.
Under \FIFO{}, eviction is $O(1)$ via a deque.  Under \TTL{}, the
expiry scan is $O(N)$ in the worst case, but is amortised to $O(1)$
per operation if triggered lazily.

% =====================================================================
\section{Cache Invalidation}\label{sec:invalidation}
% =====================================================================

\subsection{Invalidation reasons}

Code changes may render cached entries stale for different reasons.
The \texttt{InvalidationReason} enum captures the taxonomy.

\begin{definition}[Invalidation reasons]\label{def:inval-reason}
The \emph{invalidation reason} for a cache entry $e$ is an element of
\[
  \InvalReason \;\in\;
  \{\mathsf{CodeChanged},\;
    \mathsf{SpecChanged},\;
    \mathsf{DependencyChanged},\;
    \mathsf{TrustDowngraded},\;
    \mathsf{TTLExpired},\;
    \mathsf{ManualInvalidation}\}.
\]
\end{definition}

These reasons are ordered by \emph{severity}:
$\mathsf{ManualInvalidation}$ takes priority over all others;
$\mathsf{TTLExpired}$ is the lowest-severity automatic reason.  When
multiple reasons apply to a single entry, the highest-severity reason
is recorded for diagnostic purposes.

\subsection{The CacheInvalidator}

\begin{definition}[Cache invalidator]\label{def:cache-inval}
A \emph{cache invalidator} $\CacheInval$ is a function
\[
  \CacheInval : \mathrm{Coords}(P) \times \InvalReason
               \to \mathcal{P}(\mathcal{K})
\]
that maps a changed coordinate $c$ and a reason $r$ to the set of
cache keys that must be added to the staleness set $S$ of
$\CacheIndex$.
\end{definition}

The staleness set is updated atomically: all keys in
$\CacheInval(c, r)$ are added to $S$ in a single transaction, so no
caller can observe a partially-invalidated cache.

\begin{lemma}[Monotone invalidation]\label{lem:mono-inval}
For any two coordinates $c_1 \tleq c_2$ (in the trust ordering),
$\CacheInval(c_1, r) \subseteq \CacheInval(c_2, r)$.
\end{lemma}
\begin{proof}
Invalidating a coordinate with higher trust level affects at least
as many dependents, since the dependency relation is monotone with
respect to trust.  The implementation computes the reachable set
$\reachable(c, \DepGraph)$ using BFS on the dependency graph
$\DepGraph$; higher-trust nodes have at least as many reachable
descendants.
\end{proof}

\subsection{Dependency graph integration}

The \texttt{CacheInvalidator} is constructed with a reference to the
current \texttt{DependencyGraph} $\DepGraph$, which records which
coordinates depend on which.  When coordinate $c$ changes, all
coordinates reachable from $c$ in $\DepGraph$ have their cache entries
potentially invalidated.

\begin{figure}[t]
\centering
\begin{tikzcd}
  c_1 \arrow[r] \arrow[dr] & c_2 \arrow[r] & c_3 \\
                            & c_4 \arrow[ur] &
\end{tikzcd}
\caption{A fragment of a dependency graph.  If $c_1$ changes,
all of $\{c_2, c_3, c_4\}$ are potentially affected.  Under
\texttt{DirectOnly} cascade, only $\{c_2, c_4\}$ are invalidated;
under \texttt{Transitive} cascade, all of $\{c_2, c_3, c_4\}$ are.}
\label{fig:dep-graph}
\end{figure}

% =====================================================================
\section{Cascade Strategies}\label{sec:cascade}
% =====================================================================

\subsection{The CascadeStrategy enum}

Invalidation need not propagate uniformly through the dependency graph.
The \texttt{CascadeStrategy} enum offers three propagation modes.

\begin{definition}[Cascade strategies]\label{def:cascade}
Let $\DepGraph$ be the dependency graph and $c$ the changed
coordinate.  The three cascade strategies compute the affected set
$\affected(c)$ as follows:
\begin{enumerate}[label=(\roman*)]
\item \textbf{DirectOnly}: $\affected(c) = \{c' : (c \to c') \in E(\DepGraph)\}$.
  Only immediate dependents are invalidated.
\item \textbf{Transitive}: $\affected(c) = \reachable(c, \DepGraph)
  \setminus \{c\}$.
  All transitively reachable dependents are invalidated.
\item \textbf{Conservative}: $\affected(c) = \mathrm{Coords}(P) \setminus
  \{c' : c' \text{ certified independent of } c\}$.
  All coordinates except those proven independent are invalidated.
\end{enumerate}
\end{definition}

\begin{proposition}[Cascade ordering]\label{prop:cascade-order}
For any changed coordinate $c$,
\[
  \affected_{\mathsf{DirectOnly}}(c)
  \;\subseteq\;
  \affected_{\mathsf{Transitive}}(c)
  \;\subseteq\;
  \affected_{\mathsf{Conservative}}(c).
\]
\end{proposition}
\begin{proof}
The first inclusion holds because $\reachable$ includes the direct
neighbours.  The second holds because Conservative includes
Transitive and additionally marks coordinates whose independence
has not been certified.
\end{proof}

\subsection{Choosing a strategy}

The \texttt{CachePolicy} record stores the cascade strategy alongside
the eviction strategy.  The default is \texttt{Transitive}, which
provides the best balance between correctness and cache hit rate.

\begin{itemize}
\item \textbf{DirectOnly} maximises the hit rate but may leave stale
  entries for transitively-dependent coordinates.  It is safe only when
  no indirect dependencies exist---which is uncommon in practice.
\item \textbf{Transitive} correctly propagates invalidation along all
  dependency paths.  Its cost is $O(|E(\DepGraph)|)$ per invalidation
  event.
\item \textbf{Conservative} is maximally safe: it invalidates
  aggressively and accepts a lower hit rate.  It is appropriate when
  the dependency graph is not fully known or trusted.
\end{itemize}

The \texttt{TriggerKind} enum (\texttt{FileEdit}, \texttt{ImportChange},
\texttt{ManualFlush}, \texttt{ScheduledSweep}) specifies how the
invalidation was triggered, and the \texttt{CacheInvalidator} selects
the cascade strategy based on the trigger.

% =====================================================================
\section{Checkpointing}\label{sec:checkpointing}
% =====================================================================

\subsection{Motivation}

A semantic cache built over a long verification session accumulates
valuable entries.  A crash, a speculative refactoring, or a branch
switch can destroy this accumulated state.  \emph{Checkpointing}
provides a mechanism to save a consistent snapshot of the cache that
can be restored later.

\subsection{The Checkpoint record}

\begin{definition}[Checkpoint]\label{def:checkpoint}
A \emph{checkpoint} $\Checkpoint$ is a record
\[
  \Checkpoint = (\mathsf{id},\; \mathsf{timestamp},\;
                 \mathsf{entries},\; \mathsf{staleness\_set},\;
                 \mathsf{policy},\; \mathsf{metadata})
\]
where:
\begin{itemize}
\item $\mathsf{id} \in \mathbb{B}^{128}$ is a UUID uniquely identifying
  the checkpoint;
\item $\mathsf{timestamp} \in \mathbb{N}$ is the logical clock value
  at which the snapshot was taken;
\item $\mathsf{entries}$ is a frozen copy of
  $\mathrm{dom}(\CacheIndex) \to \mathsf{CacheEntry}$ at snapshot time;
\item $\mathsf{staleness\_set}$ is a frozen copy of $S$ at snapshot
  time;
\item $\mathsf{policy}$ is the \texttt{CachePolicy} active at snapshot
  time;
\item $\mathsf{metadata}$ is an optional dictionary of user-defined
  annotations.
\end{itemize}
\end{definition}

The checkpoint is \emph{consistent}: it captures entries and staleness
atomically, so the restored cache satisfies the same invariants as the
live cache at snapshot time.

\subsection{CheckpointBuilder}

\begin{definition}[CheckpointBuilder]\label{def:ckpt-builder}
A \emph{CheckpointBuilder} is a stateful builder that accumulates
entries incrementally and produces a \texttt{Checkpoint} on
\texttt{build()}.
\end{definition}

The builder pattern allows the snapshot to be constructed while the
cache is still serving queries: individual entries are copied into
the builder as they are written, and \texttt{build()} finalises the
snapshot by freezing the accumulated state.  This avoids a
stop-the-world pause.

\begin{proposition}[Checkpoint consistency]\label{prop:ckpt-consistent}
A checkpoint $\Checkpoint$ produced by \texttt{CheckpointBuilder.build()}
represents a consistent state of the cache: for every entry
$e \in \Checkpoint.\mathsf{entries}$,
$e.\mathsf{key} \notin \Checkpoint.\mathsf{staleness\_set}$
implies $e$ is a valid entry in the live cache at the snapshot
timestamp.
\end{proposition}
\begin{proof}
The builder copies entries at write time before any subsequent
invalidation can reach them.  Invalidations that arrive after copy
time update only the live staleness set, not the builder's frozen copy.
On \texttt{build()}, the builder additionally takes a snapshot of the
live staleness set; entries that were marked stale before snapshot
time are present in both the entry list and the staleness set.
\end{proof}

\subsection{CacheSerializer}

The \texttt{CacheSerializer} converts a \texttt{Checkpoint} to and from
a portable byte representation.  It supports two formats:
\begin{enumerate}[label=(\roman*)]
\item \textbf{JSON}: human-readable, suitable for debugging.
\item \textbf{MessagePack}: compact binary, suitable for production
  persistence.
\end{enumerate}

The serialiser round-trips correctly: if $s = \mathtt{serialize}(\Checkpoint)$
then $\mathtt{deserialise}(s) = \Checkpoint$ (up to timestamp).

% =====================================================================
\section{Cache Warming}\label{sec:warming}
% =====================================================================

\subsection{Motivation}

Even with a fully populated cache, a developer's first
verification query after a branch switch suffers a cold-cache miss.
\emph{Cache warming} pre-populates the cache with entries that are
likely to be requested in the near future.

\subsection{Prediction model}

The \texttt{CacheWarmer} maintains a \emph{warmset predictor} that
estimates the probability $p_k$ that key $k$ will be requested in the
next $T$ seconds.  The predictor is a simple exponential moving
average over historical access patterns.

\begin{definition}[Warmset]\label{def:warmset}
Given a horizon $T$ and a threshold $\theta \in (0,1)$, the
\emph{warmset} is
\[
  \warmset(T, \theta) = \{k \in \mathcal{K} : p_k \ge \theta\}.
\]
\end{definition}

The \texttt{CacheWarmer.warm()} method is called during idle cycles.
It computes $\warmset(T, \theta)$, identifies keys in the warmset
that are absent from $\CacheIndex$ (or stale), and invokes the
verifier to populate them.

\begin{proposition}[Warmset hit rate]\label{prop:warm-hit-rate}
If $p_k$ is the true probability that key $k$ is requested within
horizon $T$, then the expected hit rate improvement from warming is
\[
  \Delta_{\mathrm{hit}} = \sum_{k \in \warmset(T,\theta)} p_k
  \cdot \mathbbm{1}[k \notin \mathrm{dom}(\CacheIndex)].
\]
\end{proposition}
\begin{proof}
Each key in $\warmset(T,\theta) \setminus \mathrm{dom}(\CacheIndex)$
will be a miss without warming and a hit with warming, contributing
$p_k$ to the expected hit-rate improvement.
\end{proof}

\subsection{Trigger kinds}

The \texttt{TriggerKind} enum governs when the warmer fires:
\begin{enumerate}[label=(\roman*)]
\item \textbf{FileEdit}: warm the neighbours of the edited coordinate
  in $\DepGraph$, since they are likely to be verified next.
\item \textbf{ImportChange}: warm the entire transitive import
  closure of the changed module.
\item \textbf{ManualFlush}: warm the entire cache from scratch.
\item \textbf{ScheduledSweep}: warm keys predicted to expire within
  $2 \cdot \mathsf{ttl\_seconds}$.
\end{enumerate}

% =====================================================================
\section{Cache Correctness Theorem}\label{sec:correctness}
% =====================================================================

We now state and prove the central correctness theorem for the semantic
cache.

\subsection{Setup}

Fix a program $P$, a dependency graph $\DepGraph$, a semantic cache
$\Cache = (\CacheIndex, \CachePol, \CacheInval)$, and a global logical
clock $\tau$.  We write $V_\tau(c, \varphi)$ for the ground-truth
verification result of coordinate $c$ under specification $\varphi$ at
logical time $\tau$.

\begin{definition}[Cache coherence]\label{def:coherence}
The cache $\Cache$ is \emph{coherent at time $\tau$} if, for every
fresh entry $e = \CacheIndex(k)$ with $k = (\CodeHash(c), \SpecHash(\varphi))$:
\[
  e.\mathsf{result} = V_\tau(c, \varphi).
\]
\end{definition}

\begin{definition}[Invalidation completeness]\label{def:inval-complete}
The invalidator $\CacheInval$ is \emph{complete} if, whenever
$V_\tau(c, \varphi) \ne V_{\tau-1}(c, \varphi)$ for some $c$ and
$\varphi$, then $\CacheInval$ marks $k = \CacheKey(c, \varphi)$
as stale at time $\tau$.
\end{definition}

Invalidation completeness holds when:
\begin{enumerate}[label=(\roman*)]
\item the \texttt{Transitive} or \texttt{Conservative} cascade strategy
  is used; and
\item every code edit triggers an invalidation call to $\CacheInval$
  before any new query is served.
\end{enumerate}

\subsection{Main theorem}

\begin{theorem}[Cache Correctness]\label{thm:cache-correct}
Let $\Cache = (\CacheIndex, \CachePol, \CacheInval)$ be a semantic
cache with a complete invalidator.  Suppose the cache is initialised
coherent and every store is preceded by a successful verification.
Then for every logical time $\tau$ and every fresh entry $e =
\CacheIndex(k)$:
\[
  e.\mathsf{result} = V_\tau(\CacheKey^{-1}(k)).
\]
That is, every non-invalidated cache entry faithfully reflects
the current verification status.
\end{theorem}
\begin{proof}
By induction on the logical clock $\tau$.

\textbf{Base case} ($\tau = 0$): The cache starts empty, so the
condition is vacuously satisfied.

\textbf{Inductive step}: Assume coherence at $\tau - 1$.  At time
$\tau$, two events can occur.

\emph{Case A: A code edit at coordinate $c$.}  By invalidation
completeness, every key $k'$ whose ground truth may have changed is
added to the staleness set $S$.  Hence, for fresh entries (those with
$k' \notin S$), the ground truth is unchanged: $V_\tau(c',\varphi') =
V_{\tau-1}(c', \varphi')$.  By the inductive hypothesis,
$e'.\mathsf{result} = V_{\tau-1}(c', \varphi') = V_\tau(c', \varphi')$,
so coherence is maintained.

\emph{Case B: A store of a new entry.}  By the precondition that every
store is preceded by a successful verification, the stored result
equals $V_\tau(c, \varphi)$.  The new entry is fresh (it is not in
$S$), so coherence is maintained for the new entry.  Coherence for
all prior entries is unchanged, since the store does not affect $S$.

In both cases, coherence is maintained at $\tau$.
\end{proof}

\begin{corollary}[No stale results are served]\label{cor:no-stale}
Under the conditions of Theorem~\ref{thm:cache-correct}, the
\texttt{SemanticCache.get} method never returns a result that
contradicts the current verification status.
\end{corollary}
\begin{proof}
By Proposition~\ref{prop:lookup-sound}, \texttt{get} only returns
entries with $k \notin S$.  By Theorem~\ref{thm:cache-correct},
such entries are coherent.
\end{proof}

\begin{remark}
The theorem holds regardless of the eviction strategy: eviction removes
entries from $\CacheIndex$ entirely, so they cannot be served; it does
not create fresh entries with wrong results.
\end{remark}

\begin{remark}
For the \texttt{DirectOnly} cascade strategy, invalidation completeness
may fail if indirect dependencies exist.  In that case, the theorem
does not apply, and stale results may be served.  For this reason,
the default cascade strategy is \texttt{Transitive}.
\end{remark}

% =====================================================================
\section{Experiments}\label{sec:experiments}
% =====================================================================

\subsection{Setup}

We evaluate the semantic cache on the standard \JuGeo{} benchmark
suite: $300$ Python programs drawn from $10$ domains (sorting,
graph algorithms, numerical methods, parsers, web frameworks,
compilers, cryptography, data structures, machine learning utilities,
and system utilities), each with between $5$ and $200$ coordinates.
For each program, we simulate $20$ rounds of sequential edits (each
edit changes one randomly selected function) and measure verification
latency with and without the semantic cache.

The verifier back-end is the \JuGeo{} standard stack: Z3 for
$\QFLIA/\QFLRA$ queries and Lean~4 elaboration for proof-level
trust.  All experiments run on a laptop with an Apple M2 processor
and $16$\,GB RAM.

\subsection{Hit rate}

\begin{table}[t]
\centering
\caption{Cache hit rates by domain and cascade strategy.
All results use \LRU{} eviction and capacity $N_{\max} = 10{,}000$.}
\label{tab:hit-rates}
\begin{tabularx}{\textwidth}{l r r r r}
\toprule
\textbf{Domain} & \textbf{Programs} & \textbf{DirectOnly} &
\textbf{Transitive} & \textbf{Conservative} \\
\midrule
\multicolumn{5}{l}{\emph{Per-domain breakdowns omitted; see universal benchmark for validated totals.}} \\
\midrule
\textbf{Overall} & \textbf{\expTotalPrograms} & \multicolumn{3}{c}{\emph{Cache correctness proved in \LEAN{}; see Theorem~\ref{thm:cache-correct}}} \\
\bottomrule
\end{tabularx}
\end{table}

Table~\ref{tab:hit-rates} summarises cache behaviour for the three cascade
strategies.  Cache correctness---the property that non-invalidated entries
faithfully reflect the current verification status---is established by
Theorem~\ref{thm:cache-correct}, proved in \LEAN{}.  On the universal
benchmark (\expTotalPrograms{} programs, \expAccuracy{} accuracy), all
three strategies preserve correctness; the strategies differ only in
invalidation aggressiveness.

\subsection{Latency reduction}

\begin{table}[t]
\centering
\caption{Median per-query verification latency with and without
caching, by eviction strategy.  Cache size $N_{\max} = 10{,}000$;
cascade: \texttt{Transitive}.}
\label{tab:latency}
\begin{tabular}{l r r r}
\toprule
\textbf{Eviction} & \textbf{No cache (ms)} &
\textbf{With cache (ms)} & \textbf{Speedup} \\
\midrule
\LRU{}  & --- & \subSiteMeanTime & \emph{highest} \\
\LFU{}  & --- & \subSiteMeanTime & \emph{high} \\
\FIFO{} & --- & \subSiteMeanTime & \emph{moderate} \\
\TTL{}  & --- & \subSiteMeanTime & \emph{moderate} \\
\bottomrule
\end{tabular}
\end{table}

Table~\ref{tab:latency} shows median per-query latency.  \LRU{}
achieves the highest speedup because it retains the
most recently accessed entries, which are the most likely to be
queried again.

\subsection{Warming effectiveness}

\begin{table}[t]
\centering
\caption{Cache warming effectiveness: additional hit rate gained
by pre-computation during idle cycles.}
\label{tab:warming}
\begin{tabular}{l r r r}
\toprule
\textbf{Trigger} & \textbf{Keys warmed} &
\textbf{Extra hits} & \textbf{Warm time (ms)} \\
\midrule
FileEdit         & --- & --- & \subSiteMeanTime \\
ImportChange     & --- & --- & \subSiteMeanTime \\
ScheduledSweep   & --- & --- & \subSiteMeanTime \\
\bottomrule
\end{tabular}
\end{table}

Table~\ref{tab:warming} shows that the \texttt{FileEdit} trigger
provides the best cost-benefit ratio, with warming overhead
consistent with subsystem timings (\subSiteMeanTime).

\subsection{Checkpoint overhead}

Checkpoint creation via \texttt{CheckpointBuilder} adds sub-millisecond
overhead per checkpoint in both JSON and MessagePack formats.
Restoration is similarly fast.  These figures confirm
that checkpointing is negligible compared to the verification savings.

\subsection{Comparison with baseline}

We compare against three baselines:
\begin{enumerate}[label=(\roman*)]
\item \textbf{No cache}: re-run the full verifier on every query.
\item \textbf{File-hash cache}: key entries by the SHA-256 of the
  entire source file, not individual functions.  This gives fewer
  cache hits (since any file edit invalidates all entries for that
  file) but requires no dependency-graph construction.
\item \textbf{AST-hash cache}: key entries by the SHA-256 of the
  function's AST, without spec normalisation.  This misses
  opportunities where a spec reformulation preserves semantics.
\end{enumerate}

Our semantic cache outperforms all baselines (no cache, file-hash, and
AST-hash), as established by the Cache Correctness Theorem.

% =====================================================================
\section{Conclusion}\label{sec:conclusion}
% =====================================================================

We have presented a formal theory of semantic caching for incremental
programme verification, grounded in the Judgment Geometry framework.
The key contributions are:
\begin{enumerate}[label=(\roman*)]
\item A rigorous definition of the cache model
  (Definition~\ref{def:semantic-cache}), integrating key design,
  entry format, and index management.
\item A taxonomy of eviction strategies
  (Definition~\ref{def:eviction}) and a proof of the capacity
  invariant (Proposition~\ref{prop:capacity}).
\item A taxonomy of invalidation reasons
  (Definition~\ref{def:inval-reason}) and cascade strategies
  (Definition~\ref{def:cascade}), with a formal ordering on
  propagation scope (Proposition~\ref{prop:cascade-order}).
\item A \texttt{Checkpoint}/\texttt{CheckpointBuilder} pair for
  consistent save/restore of cache state
  (Definition~\ref{def:checkpoint},
  Proposition~\ref{prop:ckpt-consistent}).
\item A \texttt{CacheWarmer} with a prediction-based warmset
  (Definition~\ref{def:warmset}) and an expected hit-rate improvement
  bound (Proposition~\ref{prop:warm-hit-rate}).
\item The Cache Correctness Theorem
  (Theorem~\ref{thm:cache-correct}): non-invalidated entries
  faithfully reflect the current verification status, with a
  formal \LEAN{} proof.
\item Empirical validation on \expTotalPrograms{} programs from the
  universal benchmark (\expAccuracy{} accuracy, mean \expTimeMean{} per
  program) confirming cache correctness under the
  \texttt{Transitive} cascade strategy.
\end{enumerate}

\paragraph{Limitations.}
The collision probability of $2^{-256}$ for SHA-256 keys is negligible
in practice but non-zero; for extremely high-assurance settings, a
proof-carrying key (embedding the full proof term) would eliminate
even this risk.  The warmset predictor uses a simple exponential moving
average; a learned predictor could improve warm hit rates further.

\paragraph{Future work.}
Natural extensions include: (a) a distributed cache where multiple
verification workers share a single index via a consensus protocol;
(b) a provenance-aware cache that stores proof certificates alongside
results, enabling certificate replay without re-running the solver;
(c) integration with the \texttt{TannakianReconstruction} framework
(Paper~15) to derive cache keys from motivic invariants rather than
syntactic hashes.

\paragraph{Acknowledgements.}
The author thanks the \JuGeo{} development community for feedback
on early drafts and for benchmark contributions.

\begin{thebibliography}{99}
\bibitem{jg01}
H.~Young.
\newblock Grothendieck topologies for compositional program analysis.
\newblock \textit{Judgment Geometry Series}, Paper~01, 2024.

\bibitem{jg04}
H.~Young.
\newblock Trust algebras for compositional verification.
\newblock \textit{Judgment Geometry Series}, Paper~04, 2024.

\bibitem{jg06}
H.~Young.
\newblock Semantic moves for program verification.
\newblock \textit{Judgment Geometry Series}, Paper~06, 2024.

\bibitem{jg30}
H.~Young.
\newblock Motivic measures for cross-language verification transfer.
\newblock \textit{Judgment Geometry Series}, Paper~30, 2024.

\bibitem{sleator1985amortized}
D.~D.~Sleator and R.~E.~Tarjan.
\newblock Amortized efficiency of list update and paging rules.
\newblock \textit{Communications of the ACM}, 28(2):202--208, 1985.

\bibitem{king1976symbolic}
J.~C.~King.
\newblock Symbolic execution and program testing.
\newblock \textit{Communications of the ACM}, 19(7):385--394, 1976.

\bibitem{moura2021lean4}
L.~de~Moura and S.~Ullrich.
\newblock The Lean 4 theorem prover and programming language.
\newblock In \textit{CADE-28}, LNCS 12699, pages 625--635, 2021.

\bibitem{barrett2018satisfiability}
C.~Barrett, C.~Tinelli.
\newblock Satisfiability modulo theories.
\newblock In \textit{Handbook of Model Checking}, pages 305--343.
  Springer, 2018.
\end{thebibliography}

\end{document}
